\documentclass[tikz, border=6pt]{standalone}
\usepackage[fontset=fandol]{ctex}
\usetikzlibrary{positioning, arrows.meta, calc}

\begin{document}
\begin{tikzpicture}[
    font=\zihao{-5},
    base/.style={draw=black, line width=0.6pt, align=center, inner sep=3pt, minimum height=1.15cm},
    stage/.style={base, minimum width=2.9cm},
    data/.style={base, minimum width=3.2cm},
    proc/.style={base, minimum width=2.8cm},
    model/.style={base, fill=black!10, minimum width=2.2cm},
    head/.style={align=center, inner sep=2pt},
    arr/.style={-{Stealth[length=2.2mm, width=1.6mm]}, line width=0.6pt},
    link/.style={line width=0.6pt},
]

% 列中心
\def\cA{0}
\def\cB{3.6}
\def\cC{7.15}
\def\cD{10.5}
\def\cE{13.55}
% 行中心
\def\rA{0}
\def\rB{-1.85}
\def\rC{-3.7}
\def\rD{-5.55}
\def\rE{-7.45}

% ---------------- 表头 ----------------
\node[head] at (\cA, 1.1) {研究内容};
\node[head] at (\cB, 1.1) {数据};
\node[head] at ($(\cC, 1.1)!0.5!(\cD, 1.1)$) {关键方法};
\node[head] at (\cE, 1.1) {阶段产出};
\draw[line width=0.4pt] (-1.45, 0.8) -- (14.65, 0.8);

% ---------------- 第 2 章：CPT ----------------
\node[stage] (s1) at (\cA, \rA) {第 2 章\enspace 持续预训练\\（领域知识注入）};
\node[data]  (d1) at (\cB, \rA) {Qwen2.5-7B 基座\\电力语料 + 通用语料};
\node[proc]  (p1a) at (\cC, \rA) {困惑度\\双阈值筛选};
\node[proc]  (p1b) at (\cD, \rA) {1:1 领域-通用\\混合训练};
\node[model] (m1) at (\cE, \rA) {CPT 模型};

% ---------------- 第 3 章：SFT ----------------
\node[stage] (s2) at (\cA, \rB) {第 3 章\enspace 监督微调\\（任务能力培养）};
\node[data]  (d2) at (\cB, \rB) {专家种子、知识图谱\\七类业务场景};
\node[proc]  (p2a) at (\cC, \rB) {场景原生\\指令数据合成};
\node[proc]  (p2b) at (\cD, \rB) {两阶段微调\\+ DFT};
\node[model] (m2) at (\cE, \rB) {SFT 模型};

% ---------------- 第 4 章：RLAIF ----------------
\node[stage] (s3) at (\cA, \rC) {第 4 章\enspace 风格对齐\\（按用户角色控制风格）};
\node[data]  (d3) at (\cB, \rC) {风格标注数据\\三类用户群体};
\node[proc]  (p3a) at (\cC, \rC) {风格条件化 Prefix\\StyleRM-7B 奖励};
\node[proc]  (p3b) at (\cD, \rC) {GRPO\\+ Clip-Higher};
\node[model] (m3) at (\cE, \rC) {RLAIF 模型};

% ---------------- 第 5 章：DPO ----------------
\node[stage] (s4) at (\cA, \rD) {第 5 章\enspace 安全对齐\\（深层安全恢复）};
\node[data]  (d4) at (\cB, \rD) {Recovery / Counter\\偏好数据（四类风险）};
\node[proc]  (p4a) at (\cC, \rD) {同状态\\偏好对构造};
\node[proc]  (p4b) at (\cD, \rD) {HS-Weighting\\两阶段 DPO};
\node[model] (m4) at (\cE, \rD) {最终模型};

% ---------------- 实验验证 ----------------
\node[stage] (s5) at (\cA, \rE) {实验验证};
\node[base, minimum width=12.65cm] (ev) at ($(\cB, \rE)!0.5!(\cE, \rE) + (-0.125, 0)$)
    {领域任务：PowerQA、FaultDiag、RegulationQA、CrossDomainQA\\
     风格控制：RoleStyle-Eval；安全合规：URR、ASR、SRR、ORR；跨阶段能力保持};

% ---------------- 连线 ----------------
\foreach \i in {1,2,3,4} {
    \draw[link] (s\i.east) -- (d\i.west);
    \draw[arr] (d\i.east) -- (p\i a.west);
    \draw[arr] (p\i a.east) -- (p\i b.west);
    \draw[arr] (p\i b.east) -- (m\i.west);
}
\draw[link] (s5.east) -- (ev.west);

\draw[arr] (m1.south) -- node[right, font=\zihao{6}] {初始化} (m2.north);
\draw[arr] (m2.south) -- node[right, font=\zihao{6}] {初始化} (m3.north);
\draw[arr] (m3.south) -- node[right, font=\zihao{6}] {初始化} (m4.north);
\draw[arr] (m4.south) -- (m4.south |- ev.north);

\end{tikzpicture}
\end{document}
